\subsection{GRADED ALGEBRAS}

DEFINITION 1. Let \(\Delta\) be a commutative monoid, A a graded ring of type \(\Delta\) (II, § 11, no. 2), \((\mathbf{A}_{\lambda})_{\lambda \in \Delta}\) its graduation and E an A-algebra. A graduation \((\mathbf{E}_{\lambda})_{\lambda \in \Delta}\) of type \(\Delta\) on the additive group E is said to be compatible with the A-algebra structure on E if it is compatible both with the A-module and with the ring structure on E, in other words, if, for all \(\lambda, \mu\) in \(\Delta\) ,

\[
\mathrm{A} _ {\lambda} \mathrm{E} _ {\mu} \subset \mathrm{E} _ {\lambda + \mu}\tag{1}
\]

\[
\mathbf {E} _ {\lambda} \mathbf {E} _ {\mu} \subset \mathbf {E} _ {\lambda + \mu}\tag{2}
\]

The A-algebra E, with this graduation, is then called a graded algebra of type \(\Delta\) over the graded ring A.

When the graduation on A is trivial (that is (II, § 11, no. 1) \(\mathbf{A}_0 = \mathbf{A}, \mathbf{A}_{\lambda} = \{0\}\) for \(\lambda \neq 0\) ), condition (1) means that the \(\mathbf{E}_{\lambda}\) are sub-A-modules of E. This leads to the definition of the notion of graded algebra of type \(\Delta\) over a non-graded commutative ring A: A is given the trivial graduation of type \(\Delta\) and the above definition is applied.

When we consider graded A-algebras E with a unit element \(e\) , it will always be understood that \(e\) is of degree 0 (cf.~Exercise 1).

It follows that if an invertible element \(x \in E\) is homogeneous and of degree p, its inverse \(x^{-1}\) is homogeneous and of degree -p: it suffices to decompose \(x^{-1}\) as a sum of homogeneous elements in the relations \(x^{-1}x = xx^{-1} = e\) .

Let E and \(E'\) be two graded algebras of type \(\Delta\) over a graded ring A of type \(\Delta\) . An A-algebra homomorphism \(u: E \to E'\) is called a graded algebra homomorphism if \(u(E_{\lambda}) \subset E_{\lambda}'\) for all \(\lambda \in \Delta\) (where \((E_{\lambda})\) and \((E_{\lambda}')\) denote the respective graduations of E and \(E')\) ; where E and \(E'\) are associative and unital and u is unital, this condition means that u is a graded ring homomorphism (II, §11, no. 2).

Let E be a graded A-algebra of type N. E is identified with a graded A-algebra of type Z by writing \(E_{n} = \{0\}\) for n \textless{} 0.

Remark. Definition 1 can also be interpreted by saying that E is a graded A-module and that the A-linear mapping

\[
m \colon \mathbf {E} \otimes_ {\mathbf {A}} \mathbf {E} \rightarrow \mathbf {E}
\]

defining the multiplication on E (§ 1, no. 3), is homogeneous of degree 0 when E ⊗\_A E is given its graduation of type Δ (II, § 11, no. 5).

To define a graded A-algebra structure of type \(\Delta\) on the graded ring A, with E as underlying graded A-module, therefore amounts to defining for each ordered pair \((\lambda, \mu)\) of elements of \(\Delta\) a Z-bilinear mapping

\[
m _ {\lambda \mu}: \mathrm{E} _ {\lambda} \times \mathrm{E} _ {\mu} \rightarrow \mathrm{E} _ {\lambda + \mu}
\]

such that for every triple of indices \((\lambda, \mu, \nu)\) and for \(\alpha \in \mathbf{A}_{\lambda}\) , \(x \in \mathbf{E}_{\mu}\) , \(y \in \mathbf{E}_{\nu}\) , \(\alpha \cdot m_{\mu\nu}(x, y) = m_{\lambda + \mu, \nu}(\alpha x, y) = m_{\mu, \lambda + \nu}(x, \alpha y)\) .

Examples. (1) Let B be a graded ring of type \(\Delta\) ; if B is given its canonical \(\mathbf{Z}\) -algebra structure ( \(\S 1\) , no. 1, Example 2), B is a graded A-algebra ( \(\mathbf{Z}\) being given the trivial graduation).

\begin{enumerate}
\def\labelenumi{(\arabic{enumi})}
\setcounter{enumi}{1}
\item
  Let \(\mathbf{A}\) be a graded commutative ring of type \(\Delta\) and \(\mathbf{M}\) a graded A-module of type \(\Delta\) . Suppose that all the elements of the monoid \(\Delta\) are cancellable, which allows (II, § 11, no. 6) us to define on \(\operatorname{Homgr}_{\mathbf{A}}(\mathbf{M}, \mathbf{M}) = \operatorname{Endgr}_{\mathbf{A}}(\mathbf{M})\) a graded A-module structure of type \(\Delta\) ; as this graduation is compatible with the ring structure on \(\operatorname{Endgr}_{\mathbf{A}}(\mathbf{M})\) (II, § 11, no. 6), it defines a unital graded A-algebra structure on the A-algebra \(\operatorname{Endgr}_{\mathbf{A}}(\mathbf{M})\) .
\item
  Algebra of a magma. Let S be a magma and \(\phi: S \to \Delta\) a homomorphism. For all \(\lambda \in \Delta\) , we write \(S_{\lambda} = \phi^{-1}(\lambda)\) ; then \(S_{\lambda}S_{\mu} \subset S_{\lambda + \mu}\) . Let A be a graded commutative ring of type \(\Delta\) and \((A_{\lambda})_{\lambda \in \Delta}\) its graduation; we shall define a graded A-algebra structure on the algebra \(E = A^{(S)}\) of the magma S (§ 2, no. 6). To this end, let \(\mathbf{E}_{\lambda}\) denote the additive subgroup of \(\mathbf{E}\) generated by the elements of the form \(\alpha .s\) such that \(\alpha \in \mathbf{A}_{\mu}, s \in \mathrm{S}_{\nu}\) and \(\mu + \nu = \lambda\) . As the \(\mathrm{S}_{\lambda}\) are pairwise disjoint, \(\mathbf{E}\) is the direct sum of the \(\mathbf{A}_{\mu}\mathrm{S}_{\nu}\) and hence also the direct sum of the \(\mathbf{E}_{\lambda}\) and it is immediate that the \(\mathbf{E}_{\lambda}\) satisfy conditions (1) and (2) and therefore define on \(\mathbf{E}\) the desired graded A-algebra structure. If \(\mathbf{S}\) admits an identity element \(e\) , it may also be supposed that \(\phi(e) = 0\) . A particular case is the one where the graduation of the ring \(\mathbf{A}\) is trivial; then \(\mathbf{E}_{\lambda}\) is the sub-A-module of \(\mathbf{E}\) generated by \(\mathrm{S}_{\lambda}\) . More particularly, if we take \(\mathbf{S} = \mathbf{N}^{(I)}, \Delta = \mathbf{N}\) and \(\phi\) the mapping such that \(\phi((n_{i})) = \sum_{i \in I} n_{i}\) , the ring \(\mathbf{A}\) having the trivial graduation, a graduation is thus obtained on the polynomial algebra \(\mathbf{A}[\mathbf{X}_{i}]_{i \in I}\) , for which the degree of a homogeneous polynomial \(\neq 0\) is the total degree defined in §2, no. 9 (cf.~§6, no. 6).
\end{enumerate}

We now take S to be the free monoid Mo(B) of a set B (I, § 7, no. 2) and \(\phi\) the homomorphism \(\mathrm{Mo(B)} \to \mathbf{N}\) which associates with each word its length. Thus a graded A-algebra structure is obtained on the free associative algebra of the set B (§ 2, no. 7; cf.~§ 5, no. 5).

